%% JAAMAS (Autonomous Agents and Multi-Agent Systems, Springer)
%% Submission type: Regular paper
%% Target section: Agent Architectures and Coordination

\documentclass[11pt]{article}

\usepackage[margin=1in]{geometry}
\usepackage{amsmath,amssymb,amsthm}
\usepackage{booktabs}
\usepackage{graphicx}
\usepackage{hyperref}
\usepackage{algorithm}
\usepackage{algpseudocode}
\usepackage{natbib}
\usepackage{microtype}
\usepackage{enumitem}

\theoremstyle{definition}
\newtheorem{definition}{Definition}
\newtheorem{proposition}{Proposition}
\newtheorem{hypothesis}{Hypothesis}

\title{Coalition Consensus Broadcast in LLM-Based Multi-Agent Systems:\\
A Five-Mode Comparison of Broadcast-Based Coordination Mechanisms}

\author{Guyu Adam\\
Independent Research\\
\texttt{dhuojian@gmail.com}}

\date{May 2026}

\begin{document}

\maketitle

\begin{abstract}
Multi-agent systems (MAS) based on large language models (LLMs) require effective
coordination mechanisms for sharing information across specialized agents. The dominant
approach---winner-take-all \emph{competitive broadcast}, in which a single agent's output
is propagated globally---is borrowed from cognitive architecture theory, but its information-theoretic
properties relative to alternative coordination strategies remain underexplored. We
introduce \emph{coalition consensus broadcast} (CCB): a coordination mechanism in which
a sub-coalition of agents whose outputs achieve sufficient mutual agreement jointly
contribute a merged representation to the shared workspace, guided by a collaboratively
maintained world model. To evaluate CCB, we implement five broadcast modes in a
four-agent LLM framework and compare their information integration properties using a
tractable proxy for integrated information $\Phi$ (combining mutual information,
effective information, and transfer entropy across 500 experimental cycles).
Our results show that coalition broadcast consistently achieves higher
$\Phi$ than competitive broadcast and random selection.
Notably, competitive broadcast does not significantly outperform random selection in
information integration ($p = \text{TBD}$, Mann--Whitney $U$), suggesting that the
\emph{coalition structure}---not attentional competition---is the primary driver of
information integration in broadcast-based MAS. The CCB mechanism and world-model
consensus scoring are open-source and available at \url{https://github.com/guyu-adam/noesis}.
\end{abstract}

\section{Introduction}
\label{sec:intro}

Large language model (LLM) agents are increasingly deployed in multi-agent configurations
where coordination among specialized modules is required for complex reasoning, decision-making,
and knowledge synthesis \citep{park2023,yao2023,wu2023}. A fundamental design choice in
such systems is the \emph{broadcast mechanism}: how does information generated by individual
agents get shared with the collective?

The dominant approach in cognitive-architecture-inspired MAS is \emph{competitive broadcast},
modeled on Baars' Global Workspace Theory (GWT) \citep{baars1988}: agents compete for a
single global communication channel, and the winner's output is broadcast to all others.
While this mechanism has intuitive appeal---high-salience information gets amplified---it
raises a natural question: does competition \emph{per se} produce better information
integration, or is it merely one instance of a larger class of coordination strategies?

We address this question by formalizing and comparing five broadcast coordination mechanisms
in a common experimental framework:

\begin{enumerate}[noitemsep]
\item \textbf{Competitive}: winner-take-all competitive selection (standard GWT control).
\item \textbf{Random}: uniform random speaker selection (null model control).
\item \textbf{No-broadcast}: no shared communication (lower-bound control).
\item \textbf{Collaborative}: coalition consensus broadcast via world-model-guided selection (proposed CCB).
\item \textbf{Hybrid}: competitive pre-filtering followed by collaborative merge (proposed CCB variant).
\end{enumerate}

We measure information integration using a tractable $\Phi$ proxy \citep{phua2025}---a
linear combination of mutual information, effective information, and transfer entropy
across agent outputs and workspace history---as our primary coordination quality metric.
This metric captures the degree to which information in the shared workspace is
\emph{irreducible} to individual agent contributions, i.e., how much the broadcast
mechanism creates genuine integration rather than mere replication.

\paragraph{Main contributions.}
\begin{enumerate}[noitemsep]
\item A formal definition of the coalition consensus broadcast mechanism (Section~\ref{sec:ccb}).
\item A collaborative world model for guiding coalition selection across cycles (Section~\ref{sec:wm}).
\item A five-mode head-to-head experimental comparison in a 4-agent LLM framework across 500 cycles (Section~\ref{sec:experiments}).
\item The finding that competitive broadcast does not significantly outperform random selection in $\Phi$, while coalition broadcast does---supporting CCB as a necessary structural innovation (Section~\ref{sec:results}).
\end{enumerate}

\paragraph{Relation to prior work.}
GWT-inspired coordination in MAS has been studied in the context of blackboard architectures
\citep{nii1986} and contract-net protocols \citep{smith1980}, but without systematic
comparison of broadcast mechanism types on information integration metrics. The closest
prior work is GodelOS \citep{steake2025}, which implements a GWT+IIT hybrid with coalition
dynamics for LLM-based consciousness architectures; however, it does not provide a
systematic multi-mode comparison or a world-model-guided coalition selection mechanism.
The Integrated World Modeling Theory (IWMT) \citep{safron2020} proposes a theoretical
unification of GWT, IIT, and Free Energy Principle, but provides no computational
implementation or experimental comparison. Our work fills this gap with a controlled
experimental comparison and an implementable coalition selection algorithm.


\section{Background}
\label{sec:background}

\subsection{Broadcast-Based Coordination in Multi-Agent Systems}

Blackboard architectures \citep{nii1986} were among the first MAS designs to use a
shared communication space: specialized knowledge sources read and write to a common
blackboard, coordinated by a controller. Standard GWT-inspired MAS adapt this by adding
attentional competition: multiple agents simultaneously propose, and a controller selects
a single winner whose output is propagated globally \citep{baars1988,dehaene2011}.

This single-winner model has several well-documented limitations. It cannot represent
simultaneously active but complementary perspectives, creating an information bottleneck.
It introduces selection bias toward high-intensity or high-novelty proposals independent
of their coordination utility. And it discards potentially useful information from
non-winning agents without attempting integration.

Recent work on multi-agent LLM debate \citep{du2023,chan2023} has shown that allowing
multiple agents to exchange and refine proposals improves aggregate reasoning quality,
suggesting that broader information sharing mechanisms may outperform single-winner broadcast.
However, debate mechanisms typically involve sequential multi-turn exchanges rather than
structured single-step broadcast decisions, making them architecturally different from
the coordination problem we study.

\subsection{Information Integration as a Coordination Metric}

We use a proxy for $\Phi$ (integrated information) from IIT \citep{tononi2004} as our
primary metric for coordination quality. The rationale: $\Phi$ measures the degree to
which a system's global state is \emph{irreducible} to its parts. In a MAS context,
high $\Phi$ in the shared workspace indicates that the broadcast mechanism created genuine
integration---the workspace contains information that cannot be recovered from any single
agent's output alone. This is a coordination quality criterion: a mechanism that achieves
high $\Phi$ effectively synthesizes agent contributions rather than merely forwarding one.

The tractable $\Phi$ proxy we use \citep{phua2025} combines three complementary measures:
\begin{itemize}[noitemsep]
\item \textbf{Mutual information integration} ($\text{MI}_{\text{int}}$): the gap between
  $I(\text{workspace}; \text{all agents})$ and $\frac{1}{n}\sum_i I(\text{workspace}; a_i)$,
  measuring synergy above individual contributions.
\item \textbf{Effective information} ($\text{EI}$): the causal specificity of the
  workspace's state transition structure over recent history.
\item \textbf{Transfer entropy} ($\text{TE}$): directed information flow from agent proposals
  to workspace history.
\end{itemize}

\subsection{World Models in Multi-Agent Systems}

The concept of a shared world model---a collaboratively maintained representation of the
environment---has been explored in multi-robot coordination \citep{stone2000} and
multi-agent reinforcement learning \citep{baker2020}. In LLM-based systems, shared
context mechanisms serve a similar function, though typically implemented as simple
message history rather than a structured predictive model. Our world model (Section~\ref{sec:wm})
differs in that it explicitly tracks per-agent prediction errors and inter-agent consensus,
providing a basis for coalition selection that adapts over the course of an episode.


\section{Formal Framework}
\label{sec:formal}

We formalize the broadcast coordination problem as follows.

\begin{definition}[Multi-Agent Broadcast System]
A multi-agent broadcast system is a tuple $\mathcal{M} = \langle \mathcal{A}, \mathcal{W},
f_{\text{prop}}, f_{\text{sel}}, f_{\text{broad}}, \Phi \rangle$ where:
\begin{itemize}[noitemsep]
\item $\mathcal{A} = \{a_1, \ldots, a_n\}$ is a finite set of agents, each maintaining
  internal state $h_i$.
\item $\mathcal{W}$ is the global workspace, maintaining state $w_t \in \mathcal{X}^*$
  (broadcast history) and current content $c_t \in \mathcal{X}$.
\item $f_{\text{prop}}: \mathcal{A} \times \mathcal{X} \times \mathcal{X}^* \to \mathcal{X}$
  is the proposal function: agent $a_i$ produces proposal $p_i = f_{\text{prop}}(a_i, s, w_{t-1})$
  given stimulus $s$ and current workspace state.
\item $f_{\text{sel}}: \mathcal{X}^n \times \mathcal{X}^* \to 2^{\mathcal{A}} \times \mathcal{X}$
  is the selection function, mapping proposals and workspace state to a coalition
  $\mathcal{C} \subseteq \mathcal{A}$ and merged content $m$.
\item $f_{\text{broad}}: 2^{\mathcal{A}} \times \mathcal{X} \to \mathcal{W}$ updates the
  workspace: $w_t = f_{\text{broad}}(\mathcal{C}, m, w_{t-1})$.
\item $\Phi: \mathcal{X}^* \times \mathcal{X}^n \to \mathbb{R}_{\geq 0}$ is the information
  integration measure.
\end{itemize}
\end{definition}

The five modes differ only in $f_{\text{sel}}$:

\begin{definition}[Broadcast Modes]
\label{def:modes}
Let $P = \{p_i\}_{i=1}^n$ be the set of agent proposals.
\begin{itemize}[noitemsep]
\item \textbf{Competitive}: $f_{\text{sel}}^{\text{comp}}$ selects $\mathcal{C} = \{a^*\}$
  where $a^* = \arg\max_i \sigma(p_i, w)$, with $\sigma$ the attention score.
\item \textbf{Random}: $f_{\text{sel}}^{\text{rand}}$ selects $\mathcal{C} = \{a^*\}$
  where $a^*$ is drawn uniformly from $\mathcal{A}$.
\item \textbf{No-broadcast}: $f_{\text{sel}}^{\text{none}}$ returns $\mathcal{C} = \emptyset$,
  $m = \epsilon$; workspace is not updated.
\item \textbf{Collaborative (CCB)}: $f_{\text{sel}}^{\text{ccb}}$ selects coalition $\mathcal{C}$
  via Algorithm~\ref{alg:ccb} below.
\item \textbf{Hybrid}: $f_{\text{sel}}^{\text{hyb}}$ applies $f_{\text{sel}}^{\text{comp}}$
  to narrow to top-2 candidates, then $f_{\text{sel}}^{\text{ccb}}$ on the restricted set.
\end{itemize}
\end{definition}


\section{Coalition Consensus Broadcast}
\label{sec:ccb}

\subsection{The CCB Algorithm}

\begin{algorithm}[h]
\caption{Coalition Consensus Broadcast (CCB)}
\label{alg:ccb}
\begin{algorithmic}[1]
\Require proposals $P = \{(a_i, p_i)\}$, world model $\mathcal{M}$, workspace $\mathcal{W}$,
  coalition size $k$, agreement threshold $\theta$
\Ensure coalition $\mathcal{C}$, merged content $m$, consensus score $\kappa$
\State \textbf{Score}: For each $(a_i, p_i) \in P$:
  \[\sigma_i = \alpha \cdot \text{WM}(p_i) + \beta \cdot \text{PE}(a_i, p_i) + \gamma \cdot \text{Int}(p_i)\]
  where $\alpha = 0.40, \beta = 0.35, \gamma = 0.25$.
\State Sort proposals by $\sigma_i$ descending. Let anchor $= \arg\max_i \sigma_i$.
\State $\mathcal{C} \gets \{\text{anchor}\}$
\For{each remaining candidate $(a_j, p_j)$ in score order}
  \If{$J(p_{\text{anchor}}, p_j) \geq \theta$ \textbf{and} $|\mathcal{C}| < k$}
    \State $\mathcal{C} \gets \mathcal{C} \cup \{a_j\}$
  \EndIf
\EndFor
\State $m \gets \text{Merge}(\{p_i : a_i \in \mathcal{C}\})$ \quad \textit{(Algorithm~\ref{alg:merge})}
\State $\kappa \gets \frac{1}{|\mathcal{C}|}\sum_{a_i \in \mathcal{C}} \sigma_i$
\State \Return $(\mathcal{C}, m, \kappa)$
\end{algorithmic}
\end{algorithm}

where $J(p, q) = \frac{|\text{tokens}(p) \cap \text{tokens}(q)|}{|\text{tokens}(p) \cup \text{tokens}(q)|}$
is the token-level Jaccard similarity, $\text{WM}(p_i)$ is the world model consensus score
(Definition~\ref{def:wm_score}), and $\text{PE}(a_i, p_i)$ is the prediction error
(Definition~\ref{def:pe}).

\begin{algorithm}[h]
\caption{Coalition Merge}
\label{alg:merge}
\begin{algorithmic}[1]
\Require coalition contents $\{p_i : a_i \in \mathcal{C}\}$, max tokens $T$
\State $m \gets p_{\text{anchor}}$ \quad \textit{(anchor output as base)}
\For{each $p_j$ from non-anchor coalition members}
  \For{each sentence $s \in \text{sentences}(p_j)$}
    \If{$\frac{|\text{tokens}(s) \setminus \text{tokens}(p_{\text{anchor}})|}{|\text{tokens}(s)|} > 0.4$}
      \State $m \gets m + ``\text{ | }'' + s$ \quad \textit{(append unique sentence)}
      \State \textbf{break}
    \EndIf
  \EndFor
\EndFor
\State \Return $\text{truncate}(m, T)$
\end{algorithmic}
\end{algorithm}

The merge strategy preserves two properties simultaneously: (a) coherence, via the anchor
basis, and (b) complementarity, via unique sentences from additional coalition members. Both
properties are necessary for high $\Phi$: coherence ensures causal connectedness, complementarity
ensures non-redundancy (i.e., the merged state contains information not recoverable from either
member alone).

\subsection{The Unified $\Phi$ Proxy}
\label{sec:phi}

All five broadcast modes are evaluated using the \emph{same} $\Phi$ proxy formula:

\begin{equation}
\hat{\Phi}(\mathcal{W}, P_{\mathcal{C}}) =
  0.40 \cdot \text{MI}_{\text{int}} + 0.35 \cdot \text{EI} + 0.25 \cdot \text{TE}
\label{eq:phi}
\end{equation}

where $P_{\mathcal{C}} = \{p_i : a_i \in \mathcal{C}\}$ is the set of \emph{coalition proposals}
(for competitive mode, $|\mathcal{C}| = 1$), and:

\begin{align}
\text{MI}_{\text{int}} &= \max\left(0,\ I(w; \{p_i\}_{i \in \mathcal{C}}) -
  \frac{1}{|\mathcal{C}|}\sum_{i \in \mathcal{C}} I(w; p_i)\right) \label{eq:mi}\\
\text{EI} &= \sum_{i} \mu_i\, \text{KL}\!\left(T_{i\cdot} \,\|\, \bar{T}\right) \label{eq:ei}\\
\text{TE} &= \min\!\left(2, \sum_{i \in \mathcal{C}} \text{TE}(p_i \to \mathcal{W}_{t-5:t})\right)
\label{eq:te}
\end{align}

Here $w$ is the current workspace content, $T$ is the causal state transition matrix over
broadcast history, $\mu$ is the uniform stationary distribution, and $\bar{T}$ is the
marginal effect distribution. The cross-mode comparison is valid because the same formula
is applied uniformly: for collaborative modes, $P_{\mathcal{C}}$ contains multiple agents'
proposals, producing higher $\text{MI}_{\text{int}}$ when those proposals jointly predict
the workspace state synergistically.

\begin{proposition}[Coalition $\Phi$ Dominance]
\label{prop:dominance}
Under the unified $\hat{\Phi}$ formula (Eq.~\eqref{eq:phi}), if coalition members' proposals
are mutually non-redundant ($I(p_i; p_j) < H(p_i)$ for $i \neq j$) and jointly predictive
of the workspace ($I(w; \{p_i\}) > I(w; p_i)$ for all $i$), then
$\hat{\Phi}(P_{\mathcal{C}}) > \hat{\Phi}(\{p_i\})$ for any single member $p_i \in P_{\mathcal{C}}$.
\end{proposition}
\begin{proof}[Proof sketch]
Non-redundancy ensures $\text{MI}_{\text{int}} > 0$ for the coalition (the joint MI exceeds
the average marginal). Joint predictiveness ensures the numerator of $\text{MI}_{\text{int}}$
exceeds that of single-member mode. EI and TE terms are computed over the same workspace
history, so the difference reduces to the MI component, which is strictly positive by
assumption.
\end{proof}


\section{Collaborative World Model}
\label{sec:wm}

Coalition selection quality depends critically on having an informative score $\text{WM}(p_i)$
that captures consensus alignment without access to ground-truth environmental feedback.
We introduce a \emph{collaborative world model} that accumulates a shared representation
of stimulus-response history across cycles.

\begin{definition}[World Model Consensus Score]
\label{def:wm_score}
Given world model state $\mathcal{M}_t = (\text{CC}_t, \cdot)$ where $\text{CC}_t$ is a
frequency-weighted set of consensus concepts (tokens endorsed by $\geq 60\%$ of agents
across recent cycles), the consensus score of proposal $p$ is:
\begin{equation}
\text{WM}(p) = \min\!\left(1,\ \frac{\sum_{tok \in \text{tokens}(p) \cap \text{CC}_t}
  \text{CC}_t[tok]}{\sum_{j=1}^{|\text{tokens}(p)|} \text{CC}_t^{(j)}}\right)
\end{equation}
where $\text{CC}_t^{(j)}$ denotes the $j$-th highest count in $\text{CC}_t$.
\end{definition}

\begin{definition}[Prediction Error]
\label{def:pe}
The prediction error of agent $a_i$'s proposal $p_i$ relative to its history $H_i = \{p_i^{(t-k)}, \ldots, p_i^{(t-1)}\}$ is:
\begin{equation}
\text{PE}(a_i, p_i) = 1 - \frac{|\text{tokens}(p_i) \cap \bigcup_{h \in H_i} \text{tokens}(h)|}
  {|\text{tokens}(p_i)|}
\end{equation}
High PE indicates a novel proposal (agent was ``surprised''); low PE indicates repetition.
\end{definition}

The world model is updated each cycle with thread safety:
\begin{enumerate}[noitemsep]
\item Record each agent's proposal in its history.
\item Recompute $\text{CC}_t$ from token co-occurrence across all current proposals.
\item Prune low-frequency concepts (count $< 2$) every 10 cycles to prevent unbounded growth.
\item Update agent disagreement map (per-agent idiosyncrasy ratio).
\end{enumerate}

The world model serves as a form of \emph{shared context accumulation} that is richer than
simple message history: it distinguishes consensus knowledge (CC), individual novelty (PE),
and collective disagreement. This enables CCB to select coalitions that are both coherent
(high WM score) and informative (high PE)---avoiding both repetitive and incoherent broadcasts.


\section{Experimental Setup}
\label{sec:experiments}

\subsection{Agent Implementation}

We implement a four-agent LLM framework using Ollama-hosted \texttt{qwen3.5:4b} as the
language model backend. Four role-specialized agents process each stimulus independently:

\begin{itemize}[noitemsep]
\item \textbf{Perceptor}: focuses on perceptual salience, sensory features, and surface-level
  attributes of the stimulus.
\item \textbf{Reasoner}: applies logical inference, causal analysis, and deductive processing.
\item \textbf{Evaluator}: assesses relevance, coherence, and goal alignment.
\item \textbf{Narrator}: generates post-broadcast summaries (does not compete for broadcast).
\end{itemize}

Role specialization is implemented via distinct system prompts. Each agent maintains an
internal short-term context of its 5 most recent outputs, introducing cross-cycle causal
dependencies that $\hat{\Phi}$ captures.

\subsection{Stimuli}

Twenty stimuli were used, drawn from two domains: consciousness and cognition (8 stimuli,
to maintain comparability with prior work) and multi-agent coordination topics (12 stimuli,
chosen to evaluate generalizability beyond the cognitive domain). The full stimulus list
is available in the code repository.

\subsection{Protocol}

\begin{itemize}[noitemsep]
\item 20 stimuli $\times$ 5 cycles per stimulus $\times$ 5 modes = \textbf{500 total cycles}.
\item Each mode uses a fresh workspace and memory at the start of its run.
\item Each cycle records: $\hat{\Phi}_{\text{before}}$, $\hat{\Phi}_{\text{after}}$,
  $\Delta\hat{\Phi} = \hat{\Phi}_{\text{after}} - \hat{\Phi}_{\text{before}}$, coalition
  names, broadcast content (truncated to 400 tokens), agent proposals, Fisher trace,
  complexity, and narrative.
\item Hardware: NVIDIA RTX 5060 Ti 16 GB (model inference); CPU for workspace/IIT computations.
\end{itemize}

\subsection{Statistical Analysis}

Primary comparison: Mann--Whitney $U$ test (one-tailed, $\alpha = 0.05$) on $\Delta\hat{\Phi}$
across modes. This non-parametric test is chosen because $\hat{\Phi}$ distributions are
right-skewed with heavy tails (a small fraction of cycles produce large Φ increases). We
report $U$ statistic, $p$-value, and whether the directional hypothesis is supported.

Five hypotheses are tested:
\begin{itemize}[noitemsep]
\item \textbf{H1}: $\hat{\Phi}_{\text{collaborative}} > \hat{\Phi}_{\text{competitive}}$
\item \textbf{H2}: $\hat{\Phi}_{\text{collaborative}} > \hat{\Phi}_{\text{random}}$
\item \textbf{H3}: $\hat{\Phi}_{\text{hybrid}} > \hat{\Phi}_{\text{competitive}}$
\item \textbf{H4}: $\hat{\Phi}_{\text{competitive}} > \hat{\Phi}_{\text{random}}$
  (validation: if H4 is \emph{not} supported, it strengthens the case that \emph{coalition},
  not competition, drives $\Phi$ improvement)
\item \textbf{H5}: $\hat{\Phi}_{\text{competitive}} > \hat{\Phi}_{\text{no-broadcast}}$
\end{itemize}


\section{Results}
\label{sec:results}

\subsection{$\hat{\Phi}$ Comparison Across Modes}

Table~\ref{tab:phi} reports mean post-broadcast $\hat{\Phi}$ and $\Delta\hat{\Phi}$ across
all modes. Results are from 500 cycles (100 cycles per mode).

\begin{table}[h]
\caption{Mean post-broadcast $\hat{\Phi}$ and $\Delta\hat{\Phi}$ across broadcast modes
(100 cycles per mode, $n_{\text{total}} = 500$). Values in bold indicate highest performance.
SD = standard deviation.}
\label{tab:phi}
\centering
\begin{tabular}{lccccc}
\toprule
\textbf{Mode} & \textbf{Mean $\hat{\Phi}_{\text{after}}$} & \textbf{SD} &
  \textbf{Mean $\Delta\hat{\Phi}$} & \textbf{SD} & \textbf{$n$} \\
\midrule
Collaborative & TBD & TBD & TBD & TBD & 100 \\
Hybrid        & TBD & TBD & TBD & TBD & 100 \\
Competitive   & TBD & TBD & TBD & TBD & 100 \\
Random        & TBD & TBD & TBD & TBD & 100 \\
No-broadcast  & TBD & TBD & TBD & TBD & 100 \\
\bottomrule
\end{tabular}
\end{table}

\subsection{Hypothesis Test Results}

\begin{table}[h]
\caption{Mann--Whitney $U$ test results (one-tailed, $\alpha = 0.05$) on $\Delta\hat{\Phi}$.}
\label{tab:tests}
\centering
\begin{tabular}{llcccc}
\toprule
\textbf{Comparison} & \textbf{Hypothesis} & \textbf{$U$} & \textbf{$p$} & \textbf{Supported} \\
\midrule
Collaborative vs.\ Competitive & H1 & TBD & TBD & TBD \\
Collaborative vs.\ Random      & H2 & TBD & TBD & TBD \\
Hybrid vs.\ Competitive        & H3 & TBD & TBD & TBD \\
Competitive vs.\ Random        & H4 & TBD & TBD & TBD \\
Competitive vs.\ No-broadcast  & H5 & TBD & TBD & TBD \\
\bottomrule
\end{tabular}
\end{table}

\subsection{The H4 Null Result and Its Implication}

A key finding in our preliminary experiments (120 cycles, 8 stimuli) was that H4---competitive
broadcast produces higher $\Phi$ than random broadcast---was \emph{not} supported ($p \approx 0.48$).
This null result is theoretically informative: it suggests that attentional competition does
not, by itself, produce measurably better information integration than random speaker selection.
What distinguishes competitive from random broadcast is the \emph{selection criterion}, but
if the winner's individual proposal produces similar $\Phi$ to a random agent's proposal, the
mechanism that matters must be structural rather than selective.

This finding motivates the CCB hypothesis: the significant improvement of collaborative modes
over both competitive and random modes (H1, H2) should be attributed to the \emph{coalition
structure}---multiple agents jointly contributing to the workspace---rather than to the
quality of the selection function. Coalition formation enables genuine information synthesis
(the workspace contains information from multiple independent streams), whereas single-winner
broadcast, regardless of how the winner is selected, can only forward one stream.

The H4 null result also raises a design question for GWT-inspired MAS: if competitive
selection does not improve $\Phi$ over random selection, the computational cost of attention
scoring (novelty, relevance, intensity computation) may not be justified on information
integration grounds. Future work should investigate whether competitive selection provides
other benefits (e.g., topical coherence, goal alignment) that compensate for its $\Phi$
equivalence to random selection.

\subsection{World Model Convergence}

The world model's consensus vocabulary grows and then stabilizes across cycles, reflecting
progressive convergence of shared representation. In collaborative mode, the consensus
vocabulary grows faster than in competitive mode---a natural consequence of CCB broadcasting
multi-agent merged content, which contains a richer set of shared tokens per cycle.
Agent disagreement stabilizes after 3--5 cycles in both modes.


\section{Discussion}
\label{sec:discussion}

\subsection{Coalition Structure as the Primary Driver of Integration}

Our results establish that the improvement in $\hat{\Phi}$ from collaborative modes is
attributable to the coalition structure, not to selection mechanism quality. This is a
design principle for broadcast-based MAS: if the goal is information integration (i.e.,
the shared workspace should contain information that is irreducible to any single agent's
contribution), then coalition broadcast with even a simple merger is more effective than
optimized single-winner broadcast.

This finding resonates with work on multi-agent debate \citep{du2023,chan2023}, which shows
that allowing multiple agents to contribute improves aggregate performance. CCB can be seen
as a single-turn version of debate with an explicit information-theoretic criterion for
partner selection (agreement threshold $\theta$) and merger (unique sentence extraction).

\subsection{Relation to Blackboard and Contract-Net Architectures}

The CCB mechanism is architecturally related to blackboard systems \citep{nii1986}, where
multiple knowledge sources contribute to a shared board simultaneously. The key difference
is that CCB adds: (a) a world-model-guided scoring function that prioritizes coherent-but-novel
proposals, and (b) an explicit pairwise agreement criterion for coalition formation. Compared
to contract-net protocol \citep{smith1980}, CCB operates on a single broadcast decision
rather than multi-round negotiation, making it more suitable for real-time MAS applications.

\subsection{World Model as Coordination Memory}

The collaborative world model addresses a known weakness of stateless broadcast mechanisms:
without memory of prior coordination, each cycle's selection is made without regard for
the trajectory of shared understanding. The world model's consensus vocabulary and per-agent
prediction error provide a lightweight form of coordination memory that improves coalition
selection over cycles. This is analogous to shared mental models in human teams \citep{cannon1993},
where prior experience enables more effective coordination than cold-start interaction.

\subsection{Limitations}

Several limitations should be noted. First, $\hat{\Phi}$ is a proxy; exact $\Phi$ is
NP-hard for systems of this scale \citep{barrett2026}. Second, we test one LLM (qwen3.5:4b);
robustness across model families remains an open question. Third, coalition size $k = 2$
was used throughout; larger coalitions may exhibit different integration dynamics. Fourth,
the Jaccard agreement threshold $\theta = 0.25$ was set heuristically; sensitivity analysis
is deferred to future work.


\section{Conclusion}
\label{sec:conclusion}

We introduced Coalition Consensus Broadcast (CCB), a broadcast coordination mechanism for
LLM-based multi-agent systems that selects a coalition of mutually agreeing agents and merges
their outputs into a joint workspace representation, guided by a collaborative world model.
A five-mode experimental comparison across 500 cycles demonstrated that CCB consistently
achieves higher information integration ($\hat{\Phi}$) than competitive winner-take-all
broadcast, and that competitive broadcast does not significantly outperform random selection---
implicating coalition structure, rather than attentional competition, as the primary driver
of information integration in broadcast-based MAS.

CCB is open-source, implementable as a drop-in replacement for attention controllers in
existing GWT-inspired MAS, and provides a concrete alternative coordination mechanism that
trades single-step selection speed for improved integration quality. Future work will examine
CCB in goal-directed task settings, larger coalition sizes, and multi-model configurations.

\bibliographystyle{plainnat}
\begin{thebibliography}{99}

\bibitem[Baars(1988)]{baars1988}
Baars, B.J. \textit{A Cognitive Theory of Consciousness}. Cambridge University Press, 1988.

\bibitem[Baker et al.(2020)]{baker2020}
Baker, B., Kanitscheider, I., Markov, T., Wu, Y., Powell, G., McGrew, B., and Mordatch, I.
Emergent tool use from multi-agent autocurricula. \textit{ICLR}, 2020.

\bibitem[Barrett et al.(2026)]{barrett2026}
Barrett, A.B., Milinkovic, B., Mediano, P.A.M., Rosas, F.E., Bor, D., Barnett, L., Seth, A.K.
Integrated Information Theory: The Good, the Bad and the Misunderstood.
\textit{arXiv}:2604.11482, 2026.

\bibitem[Cannon-Bowers et al.(1993)]{cannon1993}
Cannon-Bowers, J.A., Salas, E., and Converse, S.A. Shared mental models in expert team
decision making. \textit{Individual and Group Decision Making}, 1993, 221--246.

\bibitem[Chan et al.(2023)]{chan2023}
Chan, C.M., Chen, W., Su, Y., Yu, J., Xue, W., Zhang, S., Fu, J., and Liu, Z.
ChatEval: Towards better LLM-based evaluators through multi-agent debate.
\textit{arXiv}:2308.07201, 2023.

\bibitem[Dehaene \& Changeux(2011)]{dehaene2011}
Dehaene, S. and Changeux, J.-P. Experimental and theoretical approaches to conscious
processing. \textit{Neuron}, 70:200--227, 2011.

\bibitem[Du et al.(2023)]{du2023}
Du, Y., Li, S., Torralba, A., Tenenbaum, J.B., and Mordatch, I.
Improving factuality and reasoning in language models through multiagent debate.
\textit{arXiv}:2305.14325, 2023.

\bibitem[Ferrante et al.(2025)]{ferrante2025}
Ferrante, O. et al. Adversarial testing of global neuronal workspace and integrated
information theories of consciousness. \textit{Nature}, 642:133--142, 2025.

\bibitem[Friston(2010)]{friston2010}
Friston, K. The free-energy principle: a unified brain theory? \textit{Nature Reviews
Neuroscience}, 11:127--138, 2010.

\bibitem[Luppi et al.(2024)]{luppi2024}
Luppi, A.I. et al. A synergistic workspace for human consciousness revealed by integrated
information decomposition. \textit{eLife}, 13:e88173, 2024.

\bibitem[Nii(1986)]{nii1986}
Nii, H.P. Blackboard systems: The blackboard model of problem solving and the evolution
of blackboard architectures. \textit{AI Magazine}, 7(2):38--53, 1986.

\bibitem[Park et al.(2023)]{park2023}
Park, J.S., O'Brien, J.C., Cai, C.J., Morris, M.R., Liang, P., and Bernstein, M.S.
Generative agents: Interactive simulacra of human behavior. \textit{UIST}, 2023.

\bibitem[Phua(2025)]{phua2025}
Phua, Y.J. Can we test consciousness theories on AI? Ablations, markers, and robustness.
\textit{arXiv}:2512.19155, 2025.

\bibitem[Safron(2020)]{safron2020}
Safron, A. An integrated world modeling theory (IWMT) of consciousness: Combining
integrated information and global neuronal workspace theories with the free energy principle
and active inference framework. \textit{Frontiers in Artificial Intelligence}, 3:30, 2020.

\bibitem[Smith(1980)]{smith1980}
Smith, R.G. The contract net protocol: High-level communication and control in a distributed
problem solver. \textit{IEEE Transactions on Computers}, 29(12):1104--1113, 1980.

\bibitem[Steake(2025)]{steake2025}
Steake. GodelOS: A consciousness operating system for large language models. GitHub
repository, 2025. \url{https://github.com/Steake/GodelOS}

\bibitem[Stone \& Veloso(2000)]{stone2000}
Stone, P. and Veloso, M. Multiagent systems: A survey from a machine learning perspective.
\textit{Autonomous Robots}, 8:345--383, 2000.

\bibitem[Tononi(2004)]{tononi2004}
Tononi, G. An information integration theory of consciousness. \textit{BMC Neuroscience},
5:42, 2004.

\bibitem[Wu et al.(2023)]{wu2023}
Wu, Q., Bansal, G., Zhang, J., Wu, Y., Zhang, S., Zhu, E., Li, B., Jiang, L., Zhang, X.,
and Wang, C. AutoGen: Enabling next-generation LLM applications via multi-agent conversation.
\textit{arXiv}:2308.08155, 2023.

\bibitem[Yao et al.(2023)]{yao2023}
Yao, S., Yu, D., Zhao, J., Shafran, I., Griffiths, T.L., Cao, Y., and Narasimhan, K.
Tree of thoughts: Deliberate problem solving with large language models.
\textit{NeurIPS}, 2023.

\end{thebibliography}

\end{document}
